% Part: history
% Chapter: biographies
% Section: rozsa-peter

\documentclass[../../../include/open-logic-section]{subfiles}

\begin{document}

\olfileid{his}{bio}{pet}

\olsection{R\'ozsa P\'eter}

\olphoto{peter-rozsa}{R\'ozsa P\'eter}

R\'ozsa P\'eter 原名 R\'osza Politzer，1905年2月17日出生于匈牙利布达佩斯。她因在递归函数方面的工作而闻名，这些工作对递归理论领域的建立至关重要。

P\'eter 成长于严酷的政治时期——在她十几岁时，第一次世界大战正酣——但她仍得以进入布达佩斯富庶的玛丽亚·特蕾西亚女子学校学习，并于1922年毕业。随后，她在布达佩斯的帕兹马尼·彼得大学（后更名为罗兰·厄特沃什大学）求学。起初在父亲的坚持下学习化学，后来转入数学专业，并于1927年毕业。虽然她具备教授高中数学的资格，但当时大萧条席卷全球，经济形势严峻。在此期间，P\'eter 靠做家教等零工维持生计。最终，她重返大学攻读数学研究生。她原本计划研究数论，但在发现自己的结果已经被证明后，她几乎要彻底放弃数学。在他人的鼓励下，她开始研究 Gödel 的不完全性定理，并在不知情的情况下用不同方式证明了其中的几个结果。这让她重拾信心，受 David Hilbert 基础纲领的启发，P\'eter 随后撰写了她关于递归论的首批论文。她于1935年获得博士学位，1937年成为《符号逻辑杂志》的编委。

P\'eter 的早期论文被广泛认为是递归函数论领域的奠基性贡献。在 \cite{Peter1935a} 中，她研究了不同种类的递归之间的关系。在 \cite{Peter1935b} 中，她证明了某个递归定义的函数不是原始递归的。这一结果简化了 Wilhelm Ackermann 的一个早期结果。P\'eter 简化后的函数即现今常称的 Ackermann 函数——有时更确切地称为 Ackermann--P\'eter 函数。她撰写了关于递归函数论的第一本专著 \citep{Peter1951}。

尽管她的工作重要且影响深远，P\'eter 直到1945年才获得全职教职。二战期间纳粹占领匈牙利时，由于反犹法律，P\'eter 被禁止任教。1944年，政府在布达佩斯设立了一个犹太人聚居区；该聚居区与城市其他地区隔离，由武装卫兵看守。P\'eter 被迫住在聚居区内，直至1945年该地获得解放。之后，她前往布达佩斯教师培训学院任教，自1955年起在厄特沃什·罗兰大学任教。她是匈牙利第一位获得数学学术博士头衔的女性数学家，也是第一位当选匈牙利科学院院士的女性。

P\'eter 是一位充满激情的数学教师，她更倾向于与学生一起探索数学问题的本质与美感，而非单纯授课。因此，学生们亲切地称她为“罗莎阿姨”。P\'eter 于1977年去世，享年~71岁。

\begin{reading}
更多传记资料请参见 \citep{Oconnor2014} 与 \citep{Andrasfai1986}。\citet{Tamassy1994} 对 P\'eter 进行了一次简短访谈。欲阅读关于数学的趣味读物，可参阅 P\'eter 的著作 \emph{Playing With Infinity} \citep{Peter2010}。
\end{reading}

\end{document}